\chapter{Gia Đình \& Tình Thương (Family \& Love)}

\begin{center}
  \textit{\color{truyenblue}``Family is not an important thing. It's everything.''}\\
  \textit{(Gia đình không phải là một điều quan trọng. Đó là tất cả mọi thứ.)}
\end{center}

\ngancach

\section{The Warm Dinner (Bữa Tối Ấm Áp)}
\begin{truyen}{Gia Đình}{Gắn Kết}
A father works \textbf{late}. But he always \textbf{makes} time to eat \textbf{dinner} with his \textbf{family} every night.

\textit{Một người cha làm việc muộn. Nhưng ông luôn dành thời gian ăn tối cùng gia đình mỗi tối.}

\begin{baihoc}
Bữa cơm gia đình là khoảnh khắc gắn kết tuyệt vời nhất sau một ngày dài.
\end{baihoc}
\end{truyen}

\section{The Old Blanket (Chiếc Chăn Cũ)}
\begin{truyen}{Gia Đình}{Ấm Áp}
A grandmother \textbf{knits} a blanket for her \textbf{grandson}. He \textbf{keeps} it on his bed, feeling her \textbf{love} every winter.

\textit{Một người bà đan một chiếc chăn cho cháu trai. Cậu giữ nó trên giường, cảm nhận tình yêu của bà mỗi mùa đông.}

\begin{baihoc}
Tình yêu thương gia đình mang lại sự ấm áp xua tan mọi giá lạnh.
\end{baihoc}
\end{truyen}

\section{The Broken Window (Cửa Sổ Vỡ)}
\begin{truyen}{Gia Đình}{Bao Dung}
A boy \textbf{breaks} a window. He is \textbf{scared}, but his mother \textbf{hugs} him first before \textbf{asking} what happened.

\textit{Một cậu bé làm vỡ cửa sổ. Cậu rất sợ, nhưng mẹ cậu ôm cậu trước khi hỏi chuyện gì đã xảy ra.}

\begin{baihoc}
Sự tha thứ và thấu hiểu luôn là sợi dây bền chặt của tình thân.
\end{baihoc}
\end{truyen}

\section{The Quiet Sacrifice (Sự Hy Sinh Lặng Thầm)}
\begin{truyen}{Gia Đình}{Tình Cha}
A father eats the \textbf{burnt} pieces of \textbf{toast} so his children can have the \textbf{good} ones without \textbf{knowing}.

\textit{Một người cha ăn những miếng bánh mì nướng bị cháy để các con có được những miếng ngon mà không hay biết.}

\begin{baihoc}
Tình yêu của cha mẹ thường ẩn chứa trong những hy sinh thầm lặng.
\end{baihoc}
\end{truyen}

\section{The First Bicycle (Chiếc Xe Đạp Đầu Tiên)}
\begin{truyen}{Gia Đình}{Che Chở}
A mother \textbf{runs} behind her daughter's \textbf{bicycle}, holding it \textbf{steady} until the girl can ride \textbf{alone}.

\textit{Một người mẹ chạy sau xe đạp của con gái, giữ nó thăng bằng cho đến khi cô bé có thể tự đi một mình.}

\begin{baihoc}
Gia đình luôn là chỗ dựa vững chắc nâng bước ta trên đường đời.
\end{baihoc}
\end{truyen}

\section{The Wooden Bowl (Chiếc Bát Gỗ)}
\begin{truyen}{Gia Đình}{Gương Mẫu}
An old man \textbf{drops} his food. His son gives him a \textbf{wooden} bowl. The grandson says he will do the \textbf{same} for his father. 

\textit{Một cụ ông làm rơi thức ăn. Con trai đưa cho ông một chiếc bát gỗ. Đứa cháu trai nói rằng cậu cũng sẽ làm như vậy với cha mình.}

\begin{baihoc}
Cách bạn đối xử với cha mẹ hôm nay là tấm gương cho con cái bạn ngày mai.
\end{baihoc}
\end{truyen}

\section{The Rainy Day (Ngày Mưa)}
\begin{truyen}{Gia Đình}{Niềm Vui}
It is \textbf{raining} outside. The family plays \textbf{games} inside and realizes they \textbf{do not} need sunshine to be \textbf{happy}.

\textit{Trời đang mưa bên ngoài. Cả gia đình chơi trò chơi trong nhà và nhận ra họ không cần ánh nắng mặt trời để hạnh phúc.}

\begin{baihoc}
Hạnh phúc gia đình đến từ những khoảnh khắc đơn giản bên nhau, dù ngoài kia có giông bão.
\end{baihoc}
\end{truyen}

\section{The Forgiving Brother (Khóa Chặt Căm Hờn)}
\begin{truyen}{Gia Đình}{Tình Anh Em}
Two brothers \textbf{fight} over a toy. The older brother \textbf{gives} it up, showing that his \textbf{brother} is more important.

\textit{Hai anh em tranh giành một món đồ chơi. Người anh nhường nó, cho thấy em trai mình quan trọng hơn.}

\begin{baihoc}
Trong gia đình, tình yêu thương luôn chiến thắng mọi sự nhỏ nhen ích kỷ.
\end{baihoc}
\end{truyen}

\section{The Handprint (Dấu Bàn Tay)}
\begin{truyen}{Gia Đình}{Kỷ Niệm}
A boy \textbf{puts} a messy \textbf{handprint} on the wall. His parents \textbf{frame} it instead of \textbf{cleaning} it.

\textit{Một cậu bé in một vết tay bẩn lên tường. Bố mẹ cậu đóng khung nó thay vì lau sạch.}

\begin{baihoc}
Những dấu vết tuổi thơ của con cái là tài sản quý giá nhất của cha mẹ.
\end{baihoc}
\end{truyen}

\section{The Evening Walk (Chuyến Đi Dạo)}
\begin{truyen}{Gia Đình}{Bình Dị}
An old couple \textbf{walks} slowly in the \textbf{park}. They hold \textbf{hands}, talking about their \textbf{happy} life together.

\textit{Một cặp đôi già đi dạo chầm chậm trong công viên. Họ nắm tay nhau, nói về cuộc sống hạnh phúc cùng nhau.}

\begin{baihoc}
Tình nghĩa vợ chồng bền chặt là món quà đẹp nhất của thời gian.
\end{baihoc}
\end{truyen}
